\documentclass[11pt]{amsart}
\usepackage[margin=1in]{geometry}
\usepackage{amsmath,amssymb,amsthm,booktabs,hyperref}
\hypersetup{colorlinks=true,linkcolor=blue,citecolor=blue,urlcolor=blue}
\newtheorem{theorem}{Theorem}[section]
\newtheorem{lemma}[theorem]{Lemma}
\newtheorem{corollary}[theorem]{Corollary}
\newtheorem{remark}[theorem]{Remark}
\newcommand{\C}{\mathbb C}
\newcommand{\N}{\mathbb N}
\newcommand{\Repart}{\operatorname{Re}}
\newcommand{\coeff}{[z^n]}
\title[Split-pair oscillation and a verified low rectangle]{Split-pair oscillation, the Toeplitz--Jacobi bridge,\\and a Lean-verified low zero-free rectangle}
\author{Biswajit Mondal}
\date{21 October 2026\\Unreviewed formalization preprint, version 2}
\keywords{Toeplitz minors, total positivity, entire functions, pole asymptotics, formal verification, Lean}
\begin{document}
\begin{abstract}
We give a Lean formalization of a pole-dominance criterion for sign oscillation of contiguous Toeplitz minors. A conjugate pair at the boundary of a selected set of inverse zero moduli produces two dominant terms; a finite geometric-sum argument gives both signs at arbitrarily large ranks. Repeated-pole cancellation and an ordered subset-product estimate control determinant perturbations from an analytic remainder. For an entire function, the remainder and its coefficient identity are constructed from actual simple poles of the reciprocal. A complementary determinant identity then transfers oscillation to minors of the original ordinary Taylor coefficients, with an explicit sign convention. The theorem is standalone: it assumes simple selected zeros and an isolated boundary pair, and does not replace the general total-positivity bridges used in RH reductions. Separately, a rational Taylor-model evaluator certifies nonvanishing of the actual completed xi function on the closed rectangle $0\le\Repart s\le1$, $|\Im s|\le14.1$. An attributed backport of finite-order Hadamard factorization, combined with a formal polynomial sector theorem and real polynomial approximants, gives unconditional P\'olya-frequency nonnegativity through rank 43 for the ordinary folded-xi coefficients. These finite conclusions formalize classical facts; they do not prove RH. The full sharded certificate replay and theorem audits pass for the pinned source snapshot.
\end{abstract}
\maketitle

\section{Scope and coefficient conventions}
The algebraic and analytic mechanisms here are classical. Aissen, Schoenberg and Whitney's work on totally positive sequences \cite{asw} provides the finite polynomial background. Edrei's proof already relates minors of a sequence to complementary minors of its reciprocal and uses pole information in total-positivity arguments \cite{edrei}. Toeplitz determinants and Schur functions have a broader classical and modern literature; see, for example, \cite{bumpdiaconis}. Our contribution is an explicit, checked assembly of the finite determinant expansion, its analytic perturbation bound, and the bridge for actual Taylor coefficients. We make no claim to a new total-positivity characterization.

Let $F$ be entire, with $F(0)=1$, and write
\[
 F(z)=\sum_{k\ge0}\alpha_k z^k,\qquad
 \frac1{F(z)}=\sum_{k\ge0}h_k z^k
\]
near zero. These are \emph{ordinary} coefficients:
$\alpha_k=F^{(k)}(0)/k!$. They are not the classical Jensen coefficients $k!\alpha_k$.
Set $\mu_k=(-1)^k\alpha_k$, so $\sum\mu_kz^k=F(-z)$, and extend $\mu_k=0$ for $k<0$. Define
\[
 T_{r,s}(\mu)=\det(\mu_{s+i-j})_{0\le i,j<r},\qquad
 D_s(n)=\det(h_{n+i+s-1-j})_{0\le i,j<s}.
\]
Here $r$ is the determinant size and $s$ the shift. Throughout the oscillation statements $s\ge1$. The complementary identity proved below is
\begin{equation}\label{eq:jacobi}
 D_s(n)=T_{n+s-1,s}(\mu).
\end{equation}
Thus a fixed, small reciprocal determinant controls original minors of increasing size. In the formal code the index is $n+(s-1)$, with natural-number subtraction.

The analytic theorem uses only a finite zero atlas in one closed disk. A genus-zero entire function satisfying the hypotheses is an application. The split-pair theorem does not require a general infinite canonical-product factorization. It does not allow arbitrary multiple poles or arbitrary ties at the selected modulus boundary. The separate finite-PF application below uses an actual-Xi canonical product without those genericity restrictions.

\section{Finite pole sums and oscillation}
Take distinct complex numbers $z_1,\ldots,z_m$ and nonzero residues $c_1,\ldots,c_m$. Put
\[
 u_k=\sum_{\ell=1}^m c_\ell z_\ell^k,\qquad
 M_s(n)=(u_{n+i+s-1-j})_{0\le i,j<s}.
\]
For each subset $I$ of size $s$, put
\[
 Z_I=\prod_{\ell\in I}z_\ell,\qquad
 A_I=\det\left(\sum_{\ell\in I}c_\ell z_\ell^{i+s-1-j}\right)_{i,j<s}.
\]
Multilinearity, followed by cancellation of repeated pole choices, gives the exact expansion
\begin{equation}\label{eq:subset}
 \det M_s(n)=\sum_{|I|=s} A_I Z_I^n.
\end{equation}
If $I=\{\ell_1,\ldots,\ell_s\}$, the coefficient is
\begin{equation}\label{eq:leading}
 A_I=(-1)^{s(s-1)/2}
       \left(\prod_{j=1}^s c_{\ell_j}\right)
       \prod_{i<j}(z_{\ell_j}-z_{\ell_i})^2.
\end{equation}
In particular it is nonzero. The squared Vandermonde makes this formula independent of the chosen enumeration.

\begin{lemma}[Finite dominant conjugate pair]\label{lem:finite}
Suppose conjugation permutes both the $z_\ell$ and their $c_\ell$ by the same permutation $\sigma$. Let $S$ have size $s$, let $S'=\sigma(S)\ne S$, and suppose
\[
 \Im Z_S\ne0,\qquad
 |Z_I|<|Z_S|\quad (|I|=s,\ I\notin\{S,S'\}).
\]
If a complex sequence $E_n$ satisfies $E_n/|Z_S|^n\longrightarrow0$, then
$\Repart(\det M_s(n)+E_n)$ has both signs for arbitrarily large $n$.
\end{lemma}

\begin{proof}
Conjugation gives $A_{S'}Z_{S'}^n=\overline{A_SZ_S^n}$.
After division by $|Z_S|^n$, all other terms tend to zero. The remaining expression is
$2\Repart(A_Sw^n)$, where $w=Z_S/|Z_S|$ lies on the unit circle and is nonreal.
Finite geometric sums bound the partial sums of $w^n$ and $w^{2n}$ independently of their length. Consequently the mean of $\Repart(A_Sw^n)$ tends to zero, while the mean of its square tends to $|A_S|^2/2>0$. These bounds force values of both signs bounded away from zero arbitrarily far out, and a term tending to zero cannot remove them. The Lean proof uses these elementary geometric-sum bounds; it needs no irrational-rotation theorem.
\end{proof}

This argument proves unbounded occurrences of both signs. It does not identify the first failing rank, or prove a finite-rank period formula.

\section{Analytic-error rows}
A naive entrywise estimate can repeat the largest pole modulus in several rows and lose the required decay. The proof instead cancels repeated poles before estimating mixed terms.

Assume $s\ge1$, $|S|=s$, and numbers $d>0$, $0\le\rho<d$ satisfy
\[
 |z_\ell|\ge d\quad(\ell\in S),\qquad
 |z_\ell|\le d\quad(\ell\notin S).
\]
Write $D=\prod_{\ell\in S}|z_\ell|$. If a distinct pole subset $I$ is missing $k$ entries, so $|I|+k=s$, then
\begin{equation}\label{eq:ordered}
 \left(\prod_{\ell\in I}|z_\ell|\right)d^k\le D.
\end{equation}
Indeed, divide every modulus by $d$: selected factors are at least one and all other factors at most one. This proves the comparison without requiring an ordering among moduli within $S$.

\begin{lemma}[Mixed-row determinant bound]\label{lem:mixed}
Suppose $|e_k|\le C\rho^k$, $|c_\ell|\le C$, $C\ge0$, and
$B\ge1$ bounds both $\rho$ and every $|z_\ell|$. Set
$K=CB^{2s-2}$. Then
\[
 \frac{\left|\det(u_{n+i+s-1-j}+e_{n+i+s-1-j})_{i,j<s}
                   -\det M_s(n)\right|}{D^n}
 \le (m+1)^s s!K^s\left(\frac{\rho}{d}\right)^n.
\]
\end{lemma}

\begin{proof}
Expand each row into its $m$ pole contributions and one error contribution. Subtracting the pure pole determinant leaves only choices with an error row. Two occurrences of the same pole give proportional rows, hence zero determinant, even if other rows are errors. For a surviving choice, \eqref{eq:ordered} bounds the product of distinct pole weights; the error rows contribute powers of $\rho$. At least one such row supplies the factor $\rho/d$, and additional error rows cannot increase it. The bounded offsets in the coefficient indices contribute at most $K$ per row. Leibniz's formula bounds the determinant by $s!$ times the product of its row weights, and there are at most $(m+1)^s$ choices.
\end{proof}

Since $\rho/d<1$, this establishes the normalized determinant limit required in Lemma~\ref{lem:finite}. It is a proved consequence of coefficient bounds, rather than an assumed determinant asymptotic.

\section{Entire-function assembly}
Let $a_1,\ldots,a_m$ be distinct nonzero simple zeros of $F$, closed under conjugation. Assume every zero in $|z|\le R$ is on this list, where $R>0$. The list may contain additional zeros outside that disk. Subtracting the actual principal parts constructs a remainder analytic on a neighborhood of the closed disk:
\[
 \frac1{F(z)}
   =\sum_{\ell=1}^m\frac{1/F'(a_\ell)}{z-a_\ell}+H(z)
\]
away from the poles. Removable single-pole germs are patched into one analytic remainder. Local equality near zero gives the actual coefficient identity
\[
 h_k=\sum_{\ell=1}^m c_\ell z_\ell^k+e_k,\qquad
 z_\ell=a_\ell^{-1},\qquad
 c_\ell=-\frac1{a_\ell F'(a_\ell)},\qquad
 |e_k|\le C_0R^{-k}.
\]
The Cauchy bound and this identity are proved in Lean. The residues are not freely chosen model data at this stage.

\begin{theorem}[Isolated boundary pair]\label{thm:entire}
Let $F$ be entire, $F(0)=1$, and
$F(\overline z)=\overline{F(z)}$. Suppose the zero list and disk above are given.
Let $\sigma$ be its conjugation permutation, and let $S$ have size $s$.
Assume there are indices $p\in S$ and $q\notin S$ such that
\[
 \sigma(p)=q,\quad
 \sigma(S)=(S\setminus\{p\})\cup\{q\},\quad
 |a_p|<R,\quad
 \Im\prod_{\ell\in S}a_\ell^{-1}\ne0.
\]
Assume also the boundary order and its isolated tie:
\[
 |a_y^{-1}|\le |a_x^{-1}|
       \quad(x\in S,\ y\notin S),
\]
with equality only when $x=p$ and $y=q$.
Then for every $M\in\N$ there are $r_-,r_+\ge M$ such that
\[
 T_{r_-,s}(\mu)<0,\qquad T_{r_+,s}(\mu)>0,
 \qquad \mu_k=(-1)^kF^{(k)}(0)/k!.
\]
\end{theorem}

\begin{proof}
Set $d=|a_p^{-1}|$ and $\rho=R^{-1}$; then $\rho<d$.
The boundary assumptions give the top-subset modulus conditions of Lemma~\ref{lem:mixed}. They also imply that $S$ and $\sigma(S)$ are the only size-$s$ subsets attaining the maximal product modulus. Every other subset has a strict gap, proved by the finite subset-exchange argument. Conjugation symmetry differentiates to
$F'(\overline z)=\overline{F'(z)}$, giving the required residue symmetry. Simplicity makes every residue nonzero, and distinct zeros make the Vandermonde nonzero.
The actual coefficient identity and Lemma~\ref{lem:mixed} give a normalized error tending to zero, so Lemma~\ref{lem:finite} applies to $D_s(n)$. Finally \eqref{eq:jacobi} transfers both signs to the original minors. Real-type symmetry makes these minors real. The Lean statement uses their real parts, avoiding a separate real-valued coefficient interface.
\end{proof}

\begin{corollary}[Finite products]
Theorem~\ref{thm:entire} applies to
$F(z)=\prod_{\ell=1}^m(1-z/a_\ell)$ with distinct nonzero conjugation-stable zeros, whenever the stated isolated boundary conditions hold and $R>|a_p|$ covers the needed zeros.
\end{corollary}
\begin{proof}
This polynomial is entire, has constant coefficient one, has precisely the listed simple zeros, and has the stated conjugation symmetry. The analytic theorem then applies.
\end{proof}
This corollary is a mathematical specialization of the checked theorem; the library also contains the earlier finite pole-sum theorem. We do not claim a separate new formal instantiation of arbitrary product notation in this round.

\section{The complementary determinant bridge}
For a formal series $f$ with $f_0=1$, finite lower-triangular Toeplitz matrices satisfy
$T_N(f)T_N(g)=I$ whenever $fg=1$; the boundary terms are exact ordinary convolution.
A block-coordinate proof of the classical Jacobi complementary determinant identity then gives the signed complementary-minor formula. For matrix background see \cite{hornjohnson}; \cite{bapat} states the inverse-minor identity explicitly and studies extensions. It does not require either selected minor to be invertible.

For the actual analytic functions here, the formal convolution is not an additional hypothesis. Leibniz's rule, with the identity
$\binom ni i!(n-i)!=n!$, proves that ordinary Taylor series respect multiplication. The function identity $F(z)F(z)^{-1}=1$ in a neighborhood of zero therefore gives
\[
 \operatorname{Taylor}(F)\operatorname{Taylor}(1/F)=1.
\]
Applying the signed complementary identity to $f(z)=F(-z)$ gives \eqref{eq:jacobi}. The explicit rescaling cancels the determinant sign. Omitting it would give the wrong coefficient convention.

\section{Checked declarations and reproducibility}
The project uses Lean and Mathlib v4.30.0. The main declarations in this assembly are:
\begin{center}\footnotesize
\begin{tabular}{p{.31\textwidth}p{.61\textwidth}}
\toprule
Module in \texttt{Reinmann/} & Declaration or role\\
\midrule
\texttt{SplitPairDeterminant} & \texttt{subset\_expansion}, \texttt{coefficient\_eq\_vandermonde}\\
\texttt{SplitPairMixedRows} & \texttt{coefficient\_determinant\_bound}\\
\texttt{SplitPairPatchedRemainder} & Patches the removable pole germs\\
\texttt{SplitPairCoefficientIdentity} & \texttt{actual\_reciprocal\_coefficient\_error}\\
\texttt{AnalyticSplitPair} & \texttt{entire\_split\_pair\_of\_boundary}\\
\texttt{AnalyticToeplitzJacobi} & \texttt{reciprocalMinor\_eq\_original},\newline \texttt{entire\_original\_both\_signs}\\
\bottomrule
\end{tabular}
\end{center}
\begingroup\raggedright
The checked theorem snapshot is the public repository commit
\path{7376a2336c3fe9cceaa137b247962dde3ecfc6a3}; Appendix~\ref{app:repro}
gives its toolchain, checks and release-file bindings. The 16 new declarations
in the analytic assembly are included in the snapshot's 364-declaration axiom audit.
\par\endgroup

\section{Scope limits and the motivating DH instance}
The isolated boundary hypothesis excludes general cross-boundary modulus ties and the simple-pole hypothesis excludes repeated selected zeros. These are substantive restrictions. The theorem remains separate from the repository's general PF-to-RH bridges; imposing these restrictions there would weaken their scope. No hypothesis in Theorem~\ref{thm:entire} is discharged for an unknown off-line zeta zero by this formalization.

The published computational study \emph{How sensitive are finite RH criteria? Certified Davenport--Heilbronn controls} \cite{dhcontrols} motivates this formalization. Its external interval certificates locate negative bands at shift 44 for the classical Davenport--Heilbronn control. That computational instance is distinct from a Lean instantiation of Theorem~\ref{thm:entire}: the latter would need a Lean-certified zero atlas, including coverage, simplicity, conjugation and modulus separation. Such an atlas has not been formalized. No DH numerical receipt is substituted for a Lean theorem hypothesis.

\section{The actual-Xi evaluator and low rectangle}
Use the entire normalization
\[
 \xi(s)=\frac{s(s-1)\Lambda_0(s)+1}{2},
\]
where $\Lambda_0$ is Mathlib's entire pole-subtracted completion. The extra
polynomial and constant are essential: $\Lambda_0$ itself has a different zero
set. The formal identity with the classical completed Xi is proved separately.

\begin{theorem}[Closed low rectangle]\label{thm:rectangle}
The actual function $\xi(s)$ is nonzero for
$0\le\Repart s\le1$ and $|\Im s|\le141/10$.
Consequently every zeta zero in the open critical strip satisfies
$|\Im s|\ge141/10$.
\end{theorem}

The certificate covers the quarter-rectangle
$[1/2,1]\times[0,141/10]$ by 44 rational balls.
Functional symmetry and conjugation give the full rectangle.
Each ball has a degree-12 rational complex Taylor model, with proved
coefficient error radii. Finite exponential and logarithmic enclosures,
a shifted Gamma approximation with its analytic remainder, and an
Euler--Maclaurin zeta approximation with its remainder supply the models.
The pole at $s=1$ is removed in the assembled entire function, so a tile
meeting the real edge does not assume zeta is analytic at its pole.
All bounds concern the actual analytic functions, not only formal polynomials.

For a tile centre $c$, write $a_n=\xi^{(n)}(c)/n!$. The low coefficients have
rational enclosures and all coefficients have the proved bound $|a_n|\le M$
with $M=1000$. For $0<r<1$, the derivative bound used by the certificate is
\[
 \sup_{|z-c|\le r}|\xi'(z)|\le
 \sum_{n=0}^{d-1}(n+1)|a_{n+1}|r^n+
 M r^d\left(\frac{r}{(1-r)^2}+\frac{d+1}{1-r}\right).
\]
A rational lower bound $b$ for $|\xi(c)|$, exceeding $r$ times this bound,
proves nonvanishing on the closed ball. A finite grid coverage proof checks
that every cell lies in some ball; convexity reduces coverage to its four
corners. External numerical planning chooses centres and tables, but the
Taylor-model soundness, remainders, margins and coverage are Lean propositions.
There is no appeal to an unchecked floating-point sign or to
\texttt{native\_decide}. The theorem names are
\texttt{Reinmann.XiRectangle.rectangle\_nonzero} and
\texttt{Reinmann.XiRectangle.zetaNoLowStripZero}.

\section{From the actual genus-zero product to fixed-rank PF}
Define the entire folded function $G$ by $\xi(1/2+t)=G(t^2)$, and write
$G(z)=\sum_{n\ge0}\mu_n z^n$. These $\mu_n$ are the repository's ordinary
\texttt{XiMomentCoeff}; they are not multiplied by $n!$.
PNT+'s finite-order Hadamard development \cite{pntplus} supplies the analytic
factorization machinery. The attributed backport retains 61 source files,
14,306 original lines and their Apache-2.0 authorship notices. The new
application proves exact normalization, identifies analytic multiplicities,
and derives a genus-zero product on the actual folded-Xi zero index.

The order-one growth estimate for $\xi$, with a fixed positive exponent loss,
gives order at most $3/4$ for $G$. Its Hadamard exponential factor is therefore
constant. Nonvanishing at zero identifies the constant as $G(0)>0$; inverse-zero
norm summability and the product identity are conclusions, not hypotheses.
Conjugation-stable real polynomial approximants retain repeated zeros. Their
sector positivity and locally uniform convergence transport every fixed
ordered minor to the coefficient limit. The formal polynomial sector theorem
is Schoenberg's classical sufficient criterion \cite{schoenberg,katkova}.
No fixed-rank reciprocal-PF transport is used.

\begin{theorem}[Classical fixed-rank consequence]\label{thm:pf43}
Every ordered Toeplitz minor of $(\mu_n)$ of size at most 43 is nonnegative,
at every shift.
\end{theorem}
\begin{proof}
A critical-strip zero $1/2+x+i\gamma$ has $|x|\le1/2$ and, by
Theorem~\ref{thm:rectangle}, $|\gamma|\ge14.1$. Since
$1/2\le14.1\pi/88$, its square lies outside the open sector
$|\arg z|<43\pi/44$. Schoenberg's polynomial criterion, applied to the
real approximants and passed to the limit, gives PF of order 43.
\end{proof}
The declaration is \texttt{Reinmann.XiFinitePFCertified.finitePF43}.
It has no theorem hypotheses, proves arbitrary ordered minors, and asserts
nonnegativity rather than strict positivity. This is a formalization of a
classical finite-order consequence, not a new zero-free region. The theorem
\texttt{finitePF58} still takes the named external seven-eighths proposition;
OpenAI's separate, unrefereed formalization is not imported into Lean 4.30.
Neither this application nor the low rectangle proves unbounded-rank PF,
a Green-form budget, an RH-strength ladder, or RH.

\appendix
\section{Reproducing the formal theorem snapshot}\label{app:repro}
The reproducibility pin is the full Git commit
\begin{quote}\small\ttfamily
7376a2336c3fe9cceaa137b247962dde3ecfc6a3
\end{quote}
in \url{https://github.com/sciencenerd-des/reinmann}. It uses
\texttt{leanprover/lean4:v4.30.0} and Mathlib \texttt{v4.30.0}; the exact dependency
resolution is recorded in \texttt{lake-manifest.json}. The original release at
\href{https://doi.org/10.5281/zenodo.23449490}{doi:10.5281/zenodo.23449490}
preserves that source snapshot and its 364-declaration audit. This revision uses
the later source pin below, which includes the split-pair assembly and the new
rectangle and sector programs. Manuscript files are separately checksum-bound.

From the extracted source snapshot, or from a checkout at that commit, run
\begin{quote}\small\ttfamily\raggedright
lake exe cache get\par
lake build\par
bash scripts/verify\_axiom\_clean.sh
\end{quote}
The full build passes. The source scan finds no unfinished proofs or project-declared
axioms; all 364 audited declarations depend only on \texttt{propext},
\texttt{Classical.choice}, and \texttt{Quot.sound}. The public
\href{https://github.com/sciencenerd-des/reinmann/actions/runs/39113141784}{GitHub Actions run 39113141784}
repeats the build and axiom audit for this exact commit and completed successfully.
Kernel checking establishes the stated propositions and their dependencies; it
does not certify novelty, the adequacy of mathematical hypotheses, or an unproved
numerical application.

\subsection*{Rectangle and PF43 revision snapshot}
The additional theorem-source pin is
\begin{quote}\small\ttfamily
150537bf1d94eced1c22db7660f9c32936372831.
\end{quote}
The root build imports the Hadamard and sector chain; its audit checks
2,608 theorem rows for standard axioms only. The optional numerical library is
built separately: \texttt{XiRectangle} and
\texttt{Reinmann.XiFinitePFCertified}. The hosted workflow builds shared
dependencies, 11 shards of four tiles, then assembles and audits the rectangle
and PF43. The final assembly refuses to rebuild missing tile artifacts.
All fourteen jobs in
\href{https://github.com/sciencenerd-des/reinmann/actions/runs/39470397281}{run 39470397281}
completed successfully for that exact source commit: the main build and root
audit, shared dependencies, eleven tile shards and final certificate assembly.
The final job checks the rectangle's 1,700 declarations and the unconditional
PF43 theorem separately, using only the standard foundational axioms.
The source supplement includes every numerical tile module, the complete
pinned proof source and dependency configuration, explicit audit inputs,
CI receipts and SHA256 file bindings.

The author used AI assistance for proof development, implementation and manuscript
preparation. This preprint has not completed external human peer review.

\begin{thebibliography}{12}
\bibitem{pntplus}
M. Cipollina and contributors, Hadamard-factorization development in
\emph{PrimeNumberTheoremAnd}, Apache-2.0 source,
\href{https://github.com/AlexKontorovich/PrimeNumberTheoremAnd/tree/127b50e33f6033e2147921167a117af96d9ba4ea}{commit 127b50e33f6033e2147921167a117af96d9ba4ea}, 2026.
The repository backport preserves upstream authorship and records its adaptations.
\bibitem{schoenberg}
I. J. Schoenberg, On the zeros of the generating functions of multiply positive
sequences and functions, \emph{Annals of Mathematics} \textbf{62} (1955), 447--471.
\href{https://doi.org/10.2307/1970073}{doi:10.2307/1970073}.
\bibitem{katkova}
O. M. Katkova, Multiple positivity and the Riemann zeta-function,
\emph{Computational Methods and Function Theory} \textbf{7} (2007), 13--31.
\href{https://doi.org/10.1007/BF03321628}{doi:10.1007/BF03321628}.
The preprint is \href{https://arxiv.org/abs/math/0505174}{arXiv:math/0505174} (2005).

\bibitem{asw}
M. Aissen, I. J. Schoenberg and A. M. Whitney,
On the generating functions of totally positive sequences I,
\emph{Journal d'Analyse Math\'ematique} \textbf{2} (1952), 93--103.
\href{https://doi.org/10.1007/BF02786970}{doi:10.1007/BF02786970}.
\bibitem{hornjohnson}
R. A. Horn and C. R. Johnson, \emph{Matrix Analysis}, second edition,
Cambridge University Press, 2013.
\href{https://doi.org/10.1017/CBO9781139020411}{doi:10.1017/CBO9781139020411}.
\bibitem{bapat}
R. B. Bapat, Outer inverses: Jacobi type identities and nullities of submatrices,
\emph{Linear Algebra and its Applications} \textbf{361} (2003), 107--120.
\href{https://doi.org/10.1016/S0024-3795(02)00317-8}{doi:10.1016/S0024-3795(02)00317-8}.
\bibitem{dhcontrols}
B. Mondal, \emph{How sensitive are finite RH criteria? Certified Davenport--Heilbronn controls},
unreviewed preprint, 2026.
\href{https://doi.org/10.5281/zenodo.23401370}{doi:10.5281/zenodo.23401370}.
\bibitem{edrei}
A. Edrei, Proof of a conjecture of Schoenberg on the generating function of a totally positive sequence, \emph{Canadian Journal of Mathematics} \textbf{5} (1953), 86--94.
\href{https://doi.org/10.4153/CJM-1953-010-3}{doi:10.4153/CJM-1953-010-3}.
\bibitem{bumpdiaconis}
D. Bump and P. Diaconis, Toeplitz minors, \emph{Journal of Combinatorial Theory, Series A} \textbf{97} (2002), 252--271.
\href{https://doi.org/10.1006/jcta.2001.3214}{doi:10.1006/jcta.2001.3214}.
\end{thebibliography}
\end{document}
